% GENERATED, edit the French unit or the translation table
\mxunit{example}{Example unit}
\begin{objectives}
Understand how a unit is built; use the canonical environments.
\end{objectives}
\section{A definition}
\begin{definition}[Arithmetic sequence]
A sequence $(u_n)$ is arithmetic with common difference $r$ if $u_{n+1}=u_n+r$ for every $n\in\N$.
\end{definition}
\begin{example}
With $u_0=2$ and $r=\num{0,5}$, we get $u_3=\num{3,5}$. \mxchip{5 min}
\end{example}
\begin{pitfall}
The common difference is read from the difference $u_{n+1}-u_n$, not from the quotient.
\end{pitfall}
\begin{theorem}
For every $n\in\N$, $u_n=u_0+nr$.
\end{theorem}
\begin{proof}
By induction on $n$: the equality holds for $n=0$, and if $u_n=u_0+nr$ then $u_{n+1}=u_0+(n+1)r$.
\end{proof}
\begin{exercise}[Computing terms]
Compute $u_{10}$ for $u_0=-3$ and $r=2$. Answer: $u_{10}=\mxblank{17}$.
\mxlines{2}
\end{exercise}
\begin{solution}
$u_{10}=-3+10\times2=17$.
\end{solution}
\begin{figurebox}[The terms lie on a line.]
\begin{tikzpicture}\begin{axis}[mxaxis,width=7cm,height=4cm,xmin=0,xmax=5,ymin=0,ymax=5]
\addplot[mxcurve,only marks] coordinates {(0,2)(1,2.5)(2,3)(3,3.5)(4,4)};
\end{axis}\end{tikzpicture}
\end{figurebox}
\begin{teachernote}
On \mxdate{2026-09-01}, have students restate the definition aloud.
\end{teachernote}
